\section{What a lookup costs now}
\label{sec:ch04-what-a-lookup-costs}
\index{Entity Framework Core|(}

Mosaic's docker-compose file has said the same thing since
chapter~\ref{ch:hotchocolate-quickly}: it is the spine, and later chapters
attach things to it. This is the chapter that attaches the first one.

\begin{minted}[fontsize=\footnotesize]{yaml}
  mosaic-db:
    image: postgres:18-alpine
    environment:
      POSTGRES_DB: mosaic
      POSTGRES_USER: mosaic
      POSTGRES_PASSWORD: mosaic
    volumes:
      - mosaic-db-data:/var/lib/postgresql
    healthcheck:
      test: ["CMD-SHELL", "pg_isready -U mosaic -d mosaic"]
      interval: 5s
      timeout: 3s
      retries: 10
      ...
\end{minted}

\index{PostgreSQL!container}%
The health check earns its place. Starting a fresh volume runs \code{initdb},
which brings the server up, configures it and restarts it, and a client that
connects during that window gets refused. The API service depends on this one
\code{condition: service\_healthy}, so compose holds it back until
\code{pg\_isready} answers.

The volume line cost me twenty minutes and is worth passing on. It mounts one
directory higher than it used to, and not at the \code{data} subdirectory
underneath, because PostgreSQL 18's official image moved the cluster into a
version-named directory so that \code{pg\_upgrade} can see an old one and a new
one without crossing a mount point~\autocite{dockerlibrary2026postgres18}. Mount
the old path and the container refuses to start, in a restart loop, with an
explanation in its log that is unusually good and that you will not read until
you think to look.

\subsection*{Six domains, one context, no navigations}

\index{Entity Framework Core!DbContext}%
\code{MosaicDbContext} holds a \code{DbSet} for each of the six domains and
nothing else. The mapping lives in the domains: every folder carries its own
\code{IEntityTypeConfiguration<T>}, and the context only collects them.

\begin{minted}{csharp}
protected override void OnModelCreating(ModelBuilder modelBuilder)
{
    modelBuilder.ApplyConfigurationsFromAssembly(
        typeof(MosaicDbContext).Assembly);

    UseSnakeCaseNames(modelBuilder);
}
\end{minted}

That is the same rule the resolvers follow, applied to tables. A domain's
storage is described in the folder that owns the domain, so when Catalog leaves
this service in chapter~\ref{ch:the-first-cut-extracting-catalog} its table
definition goes with it.

\index{Entity Framework Core!navigations}%
The rule with teeth is the one about navigations, and it is a rule about what
the model does \emph{not} contain. No entity here holds a navigation property to
an entity another domain owns. Reviews stores a \code{Guid} for the product, not
a \code{Product}. Ordering stores a \code{Guid} for the customer, not a
\code{Customer}. There is not a single foreign key across a domain boundary,
which means there is not a single \code{Include} to unpick later, because there
never was one.

This costs something now. A join that the database could do for free becomes two
statements and a dictionary lookup in C\#. I am paying that on purpose, and it is
worth being explicit about why rather than dressing it up: within one service the
join is cheaper, and the moment the two domains are two services with two
databases the join is not available at any price. Building the monolith as though
the boundary were already there is the only way the extraction later is a move
rather than a rewrite.

\subsection*{Which scope the context comes from}

\index{dependency injection!DbContext}%
Chapter~\ref{ch:the-life-of-a-request} read ten lines of \code{ResolverTask} and
found that a query field gets a dependency injection scope of its own, one per
resolver invocation, while a mutation field shares the request's. I said at the
time that this stops being trivia in this chapter. Here is where.

\index{Entity Framework Core!threading}%
Entity Framework Core does not support two operations running on one
\code{DbContext} at once. Query resolvers run concurrently. Put those two
sentences together and the per-resolver scope stops looking like a curiosity and
starts looking like the thing standing between you and an intermittent
\enquote{A second operation started on this context before a previous operation
completed}. The v16 documentation says as much, and warns that changing the
default scope for queries will likely produce exactly that
error~\autocite{chillicream2026efcore}.

Mosaic registers a pooled factory and a scoped context resolved from it:

\begin{minted}[fontsize=\footnotesize]{csharp}
services.AddPooledDbContextFactory<MosaicDbContext>((serviceProvider, options) =>
    options
        .UseNpgsql(connectionString)
        .AddInterceptors(
            new SqlCommandCounter(
                serviceProvider.GetRequiredService<IHttpContextAccessor>())));

services.AddScoped(serviceProvider =>
    serviceProvider.GetRequiredService<IDbContextFactory<MosaicDbContext>>()
        .CreateDbContext());
\end{minted}

\index{Entity Framework Core!RegisterDbContextFactory}%
One line in \code{Program.cs} goes with it, and it does less than its name
suggests:

\begin{minted}{csharp}
builder.AddGraphQL()
    .AddMosaic()
    .RegisterDbContextFactory<MosaicDbContext>()
\end{minted}

That call adds a parameter expression builder so the resolver compiler knows how
to hand a resolver a \code{MosaicDbContext}. It does not register the factory. I
have read the whole file it lives in, which is fifty-four lines including the
documentation comments and contains exactly two methods, so I can say that with
some confidence; the documentation says the same thing in
bold~\autocite{chillicream2026efcore}. Forget \code{AddPooledDbContextFactory}
and this line still compiles.

\subsection*{Counting what actually reaches the database}

\index{diagnostics!DbCommandInterceptor}%
Chapter~\ref{ch:hotchocolate-quickly}'s counter counts questions asked of a
domain service. That was a fair proxy for work while every question was answered
from memory, and it stops being one the moment a service can answer several
questions with one statement. So this chapter adds a second counter, and puts it
as low as it can go:

\begin{minted}[fontsize=\footnotesize]{csharp}
public sealed class SqlCommandCounter(IHttpContextAccessor httpContextAccessor)
    : DbCommandInterceptor
{
    public override ValueTask<InterceptionResult<DbDataReader>> ReaderExecutingAsync(
        DbCommand command,
        CommandEventData eventData,
        InterceptionResult<DbDataReader> result,
        CancellationToken cancellationToken = default)
    {
        Count();
        return ValueTask.FromResult(result);
    }

    ...

    private void Count()
        => httpContextAccessor.HttpContext?.RequestServices
            .GetService<RequestTimeline>()
            ?.CountSqlCommand();
}
\end{minted}

The elision hides three more overrides that do the same thing for the
synchronous reader and for the two scalar hooks, so a \code{count(*)} is
counted alongside a \code{SELECT}. An interceptor on the command is not an
estimate of round trips, it is the round trips. The awkward part is
attribution. The interceptor is attached to the
pooled factory's options, which are built once for the process, so it cannot
hold a per-request total and has to go and find one. It finds it the way
chapter~\ref{ch:the-life-of-a-request}'s listener does, through
\code{RequestServices}, reached here through the \code{IHttpContextAccessor}.

I did not assume that would work. The accessor keeps the context in an
\code{AsyncLocal}, and a DataLoader batch is dispatched from a coordinator loop
rather than from the resolver that asked for the data, which is exactly the
shape of thing that loses an ambient value. It survives, because HotChocolate
registers its batch dispatcher per request scope and starts the loop inside the
request. That is a claim I checked twice: once by reading the registration, and
once by firing ten copies of the same query at the service at the same moment
and confirming that all ten reported their own counts rather than one report of
ten times the work.

\subsection*{The number, with a price on it}

\index{N+1 problem!measured}%
The timeline line grows a field, and the catalog-with-reviews query says this:

\begin{minted}[fontsize=\footnotesize]{text}
parse - validate - compile - coerce - execute 12.674ms total 12.716ms
    (document cache hit, operation cache hit, 146 resolvers, 146 SQL)
\end{minted}

A hundred and forty-six resolvers, a hundred and forty-six lookups, a hundred
and forty-six statements. The three numbers agree because every resolver asks
one question and every question is one query, and the agreement is what is
wrong, put as plainly as it can be put.

I warmed the process with twenty requests before believing any of this, for the
reason chapter~\ref{ch:the-life-of-a-request} gave: a first request to a fresh
Mosaic reports around 205~ms and almost all of it is the runtime compiling
itself. Warm, nine of ten runs came in between 9.2 and 14.2~ms and one landed at
44.6, which is a laptop being a laptop. Against three milliseconds from a list in
memory, the same query now costs somewhere near twelve.

Read that gap carefully, because the interesting thing about it is how small it
is. PostgreSQL is on the same machine, over a loopback socket, holding
twenty-five products and a hundred and twenty reviews, with every row in cache.
This is the friendliest database a query will ever have. Move it to another host
and the per-statement cost goes up by a network round trip; give it a table
worth having an index on and it goes up again. Twelve milliseconds is the floor.
What the N+1 costs in production is that floor multiplied by conditions nobody
sets on purpose.

Before the fix, though, there is a bill to settle with the mapper.
\index{Entity Framework Core|)}
